% ===========================================================================
% Poste 19 — Impôt sur le revenu (Québec + fédéral)
% ===========================================================================
\poste{19}{Impôt sur le revenu des particuliers}{QC\_impot, CA\_impot}

\pdesc Les deux impôts sur le revenu --- québécois et fédéral --- sont
calculés par adulte selon la mécanique classique : revenu imposable, barème
progressif, crédits non remboursables, puis, au fédéral, l'abattement du
Québec. Les couples ajoutent des mécanismes de partage (montant pour
conjoint, transferts de crédits).

\pobj Financer les services publics selon la capacité de payer, tout en
reconnaissant par les crédits la situation personnelle (besoins
essentiels, âge, retraite, personne vivant seule, charges familiales).

\subsubsection*{Impôt du Québec}

\paragraph{Revenu imposable et revenu net.} Pour un adulte de revenu $r$
(travail ou pension) et d'âge $a$ :
\begin{align*}
\text{revenu imposable} &= r + PSV_{imp} - D_{trav} - RRQ_{suppl} \\
\text{revenu net} &= r + PSV_{tot} - D_{trav} - RRQ_{suppl}
\end{align*}
où $PSV_{imp}$ est la pension de la Sécurité de la vieillesse imposable et
$PSV_{tot}$ l'ensemble des prestations du poste~17 (SRG et Allocation
compris : non imposables, mais inclus dans le revenu net --- encadré de la
section~\ref{sec:bases}). $D_{trav}$ est la \emph{déduction pour
travailleur} : \pc{6} du revenu de travail, plafonnée (\D{1420} en 2025,
\D{1450} en 2026) ; nulle pour un retraité. $RRQ_{suppl}$ est la cotisation
RRQ supplémentaire, déductible (poste~1).

\paragraph{Impôt brut.} Barème progressif à quatre paliers appliqué au
revenu imposable :

\begin{center}
\small
\begin{tabular}{r r c}
\toprule
Tranche 2025 & Tranche 2026 & Taux \\
\midrule
jusqu'à \D{53255}  & jusqu'à \D{54345}  & \pc{14} \\
jusqu'à \D{106495} & jusqu'à \D{108680} & \pc{19} \\
jusqu'à \D{129590} & jusqu'à \D{132245} & \pc{24} \\
au-delà            & au-delà            & \pc{25.75} \\
\bottomrule
\end{tabular}
\end{center}

\paragraph{Crédits non remboursables.} Tous au taux du premier palier
(\pc{14}). Le montant personnel de base $BPA$ s'applique à chacun ; s'y
ajoute un \emph{montant combiné} --- personne vivant seule, âge (65~ans et
plus) et revenus de retraite --- réduit de \pc{18.75} du revenu net
excédant un seuil :
\begin{align*}
\text{combiné} &= \mathrm{(seul)} \times M_{seul}
 + \mathrm{(65+)} \times M_{âge}
 + \mathrm{(retraité)} \times \min(r,\, M_{pension}) \\
\text{combiné réduit} &= \pos{\text{combiné}
 - 0{,}1875 \times \pos{\text{revenu net} - s}} \\
\text{impôt} &= \pos{\text{impôt brut}
 - 0{,}14 \times (BPA + \text{combiné réduit})}
\end{align*}
Les cotisations sociales ne donnent pas de crédit au Québec : la déduction
pour travailleur en tient lieu. Le montant \og personne vivant seule \fg{}
s'applique aux ménages à un adulte (la famille monoparentale y a droit).

\paragraph{Couples (Québec).} Trois mécanismes s'ajoutent :
\begin{enumerate}
  \item l'adulte~1 réclame le montant combiné \emph{des deux} conjoints
        (âge et pension), réduit de \pc{18.75} du \emph{revenu net
        familial} (somme des deux revenus nets, prestations non imposables
        comprises) excédant le seuil ;
  \item un crédit \emph{non remboursable} pour frais médicaux du couple :
        $\pos{\text{prime RAMQ} - 0{,}03 \times \text{revenu net
        familial}} \times \pc{20}$ ;
  \item le \emph{transfert entre conjoints} : la partie inutilisée des
        crédits de l'un (quand ses crédits dépassent son impôt brut)
        réduit l'impôt de l'autre.
\end{enumerate}

\subsubsection*{Impôt fédéral}

\paragraph{Revenu imposable et revenu net.}
\begin{align*}
\text{revenu imposable} &= r + PSV_{imp} - RRQ_{suppl} - D_{garde} \\
\text{revenu net} &= \text{revenu imposable} + SV_{non\,imp}
\end{align*}
où $D_{garde}$ est la déduction fédérale pour frais de garde (poste~6,
réclamée par le conjoint au revenu le plus faible) et $SV_{non\,imp}$ les
prestations non imposables (SRG, Allocation) de l'adulte. Il n'y a pas de
déduction pour travailleur au fédéral.

\paragraph{Impôt brut.} Barème progressif à cinq paliers :

\begin{center}
\small
\begin{tabular}{r r c c}
\toprule
Tranche 2025 & Tranche 2026 & Taux 2025 & Taux 2026 \\
\midrule
jusqu'à \D{57375}  & jusqu'à \D{58523}  & \pc{14.5}\footnotemark & \pc{14} \\
jusqu'à \D{114750} & jusqu'à \D{117045} & \pc{20.5} & \pc{20.5} \\
jusqu'à \D{177882} & jusqu'à \D{181440} & \pc{26}   & \pc{26} \\
jusqu'à \D{253414} & jusqu'à \D{258482} & \pc{29}   & \pc{29} \\
au-delà            & au-delà            & \pc{33}   & \pc{33} \\
\bottomrule
\end{tabular}
\end{center}
\footnotetext{Taux de transition : la baisse du taux du premier palier de
15~\% à 14~\% est entrée en vigueur le 1\ier{}~juillet 2025, d'où un taux
effectif de 14,5~\% pour l'année 2025 entière.}

\paragraph{Crédits non remboursables.} Tous au taux du premier palier.
\begin{itemize}
  \item \emph{Montant personnel de base} $BPA(\cdot)$ : une partie fixe
        plus une bonification qui s'amenuise linéairement entre deux seuils
        de revenu net élevés (\D{177882} et \D{253414} en 2025) ;
  \item \emph{cotisations} : RRQ de base, RQAP et assurance-emploi
        (postes~1 à 3) ;
  \item \emph{montant canadien pour emploi} : $\min(\D{1471},\, r)$ en
        2025 (\D{1501} en 2026), travailleurs seulement ;
  \item \emph{montant en raison de l'âge} (65~ans et plus) :
        $\pos{M_{âge} - 0{,}15 \times \pos{\text{revenu net} - s_{âge}}}$,
        avec $M_{âge} = \D{9028}$ et $s_{âge} = \D{45522}$ en 2025
        (\D{9208} et \D{46432} en 2026) ;
  \item \emph{montant pour revenu de pension} : $\min(\D{2000},\, r)$ pour
        un retraité ;
  \item \emph{montant pour personne à charge admissible} (famille
        monoparentale) : un second $BPA$, même mécanique de bonification.
\end{itemize}

\paragraph{Abattement du Québec.} L'impôt fédéral net est réduit de
\pc{16.5} pour les résidents du Québec (art.~120(2) LIR) :
$\text{impôt final} = \text{impôt net} \times 0{,}835$.

\paragraph{Couples (fédéral).} Trois mécanismes :
\begin{enumerate}
  \item \emph{montant pour conjoint} : $\pos{BPA - \text{revenu net du
        conjoint}}$ --- le revenu net du conjoint \emph{inclut} ses
        prestations non imposables (fiche 110113) ;
  \item \emph{transfert} de la partie inutilisée des montants d'âge et de
        pension vers l'autre conjoint (le transfert le plus avantageux
        l'emporte) ;
  \item \emph{crédit médical} du couple :
        $\pos{\text{prime RAMQ} - 0{,}03 \times \text{revenu net}}$,
        réclamé par le conjoint qui en tire le plus, plafonné à son impôt
        résiduel.
\end{enumerate}

\pparams Récapitulatif des montants de crédits :

\begin{center}
\small
\begin{tabular}{l r r}
\toprule
Paramètre & 2025 & 2026 \\
\midrule
\multicolumn{3}{l}{\emph{Québec}} \\
Montant personnel de base ($BPA$)        & \D{18571} & \D{18952} \\
Déduction pour travailleur --- taux      & \pc{6}    & \pc{6} \\
Déduction pour travailleur --- plafond   & \D{1420}  & \D{1450} \\
Montant personne vivant seule ($M_{seul}$) & \D{2128} & \D{2172} \\
Montant en raison de l'âge ($M_{âge}$)   & \D{3906}  & \D{3986} \\
Montant pour revenus de retraite ($M_{pension}$) & \D{3470} & \D{3541} \\
Réduction du montant combiné --- taux    & \pc{18.75} & \pc{18.75} \\
Réduction du montant combiné --- seuil   & \D{42090} & \D{42955} \\
Crédit médical (couples) --- taux        & \pc{20}   & \pc{20} \\
\midrule
\multicolumn{3}{l}{\emph{Fédéral}} \\
$BPA$ --- partie de base                 & \D{14538} & \D{14829} \\
$BPA$ --- bonification                   & \D{1591}  & \D{1623} \\
Montant en raison de l'âge ($M_{âge}$)   & \D{9028}  & \D{9208} \\
Seuil de réduction du montant d'âge      & \D{45522} & \D{46432} \\
Montant pour revenu de pension (max)     & \D{2000}  & \D{2000} \\
Montant canadien pour emploi (max)       & \D{1471}  & \D{1501} \\
Abattement du Québec                     & \pc{16.5} & \pc{16.5} \\
\bottomrule
\end{tabular}
\end{center}

\prefs Québec : Loi sur les impôts (RLRQ, c.~I-3) ; Revenu Québec.
Fédéral : Loi de l'impôt sur le revenu (L.R.C. 1985, ch.~1
(5\ieme{}~suppl.)), art.~117 (barème), 118 et suivants (crédits), 120(2)
(abattement) ; Agence du revenu du Canada.
